% Updated Test Cases LaTeX Documentation
% Based on actual implemented and passing test suite (34 test cases)

\subsection{Overview}

During the development process, the Tutor Support System was thoroughly tested through both automated and manual testing approaches. The testing strategy focused on verifying core functionalities, ensuring system stability, and validating non-functional requirements such as performance, security, and scalability.

To maintain clarity and brevity in this report, we present only the most representative test cases that cover critical user scenarios and system behaviors. These selected test cases demonstrate the system's capability to handle essential workflows across all major modules. In practice, dozens of additional detailed test cases were executed during development to validate edge cases, error handling, data validation, and integration points. Those comprehensive test suites are maintained in the project repository but are not included here to avoid excessive documentation length.

\subsection{Functional Test Cases}

% tăng khoảng cách hàng cho cả section
\renewcommand{\arraystretch}{1.2}

% ---------- Functional longtable ----------
\refstepcounter{table}\label{tab:functional-testcases}
\begin{center}
    
\begin{longtable}
{p{0.06\textwidth}L{0.20\textwidth}L{0.30\textwidth}L{0.30\textwidth}p{0.05\textwidth}}
\toprule
\textbf{TC ID} & \textbf{Module / UC} & \textbf{Scenario} & \textbf{Expected result} & \textbf{Result} \\
\midrule
\endfirsthead

\toprule
\textbf{TC ID} & \textbf{Module / UC} & \textbf{Scenario} & \textbf{Expected result} & \textbf{Result} \\
\midrule
\endhead

\midrule
\multicolumn{5}{r}{\textit{Continued on next page}}\\
\bottomrule
\endfoot

\bottomrule
\caption*{\textbf{Table \thetable}. Functional test cases for core Tutor Support System features.}
\endlastfoot

F-01 & Account Management (Login) &
The user accesses the system, clicks ``Login with HCMUT\_SSO'', and successfully authenticates on the HCMUT\_SSO portal. &
The system receives a valid assertion from HCMUT\_SSO, creates an application session, and redirects the user to the dashboard with the correct role (Student/Tutor/Coordinator/Admin). &
Pass \\ \midrule

F-02 & Account Management (Profile) &
The user is logged in, opens the Profile page, updates local fields (bio, phone number, preferences), and saves the form using PUT request to /users/profile endpoint. &
Profile information is persisted to the database. When the page is reloaded, the updated data is shown correctly. Fields synchronized from HCMUT\_DATACORE are read-only and cannot be edited. &
Pass \\ \midrule

F-03 & Account Management (Session Validation) &
A logged-in user performs multiple actions (navigate between pages, call APIs) while their session token remains valid, then waits until the token expires. &
All authenticated requests succeed while the token is valid. After expiration, the system returns 401 Unauthorized and redirects to the login page without exposing protected data. &
Pass \\ \midrule

F-04 & Account Management (Logout) &
A user clicks the Logout button while having an active session. &
The system clears the session token from client storage, terminates the server-side session (if applicable), and redirects the user to the login page. Subsequent attempts to access protected routes fail with 401. &
Pass \\ \midrule

F-05 & Registration -- Student registration (UC-REG-01) &
A student logs in, opens the Registration form via /students/register endpoint, submits their university ID to register as a student in the system. &
A new student profile is created and linked to the user account. The student can now proceed to register for courses and participate in tutoring sessions. &
Pass \\ \midrule

F-06 & Registration -- Tutor subject registration (UC-REG-01) &
A tutor logs in, opens the Registration form via /tutors/register-subject endpoint, selects the course they want to tutor, specifies their expertise level, and submits. &
A new tutor subject registration with specified expertise level is created in state \textit{Pending}. It is visible to the coordinator in the list of registrations via /coordinator/tutor-registrations that require review. &
Pass \\ \midrule

F-07 & Registration -- View registration result (UC-REG-02) &
A tutor who has submitted a subject registration opens the Registration Status page via /tutors/my-registrations to check the status of their pending registrations. &
The system displays all subject registrations submitted by the tutor with accurate statuses (\textit{Pending}, \textit{Approved}, \textit{Rejected}). If rejected, the reason provided by the coordinator is visible. &
Pass \\ \midrule

F-08 & Registration review (UC-REG-03) &
A coordinator logs in, opens the list of \textit{Pending} tutor subject registrations via /coordinator/tutor-registrations, inspects the details of a tutor registration, and decides to approve or reject it. &
The registration state is updated correctly to \textit{Approved} or \textit{Rejected}. The corresponding tutor is notified. Only \textit{Approved} registrations are usable for scheduling. &
Pass \\ \midrule

F-09 & Scheduling -- Tutor creates session (UC-SCH-01) &
A tutor with at least one \textit{Approved} subject registration opens the Scheduling module via /sessions/ endpoint, creates a new session by specifying subject, scheduled date, duration, and maximum number of students. &
New sessions are created in state \textit{Proposed}, linked to the correct tutor and course. The session parameters including max students and duration are stored, while the student participant list is not yet finalized. &
Pass \\ \midrule

F-10 & Scheduling -- Set available time &
A tutor accesses their schedule preferences via /schedule-preferences/ endpoint, defines or updates recurring weekly time slots (e.g., Monday 2-4 PM, Wednesday 9-11 AM) for when they are available to tutor. &
The availability information is persisted as schedule preferences and becomes selectable when creating sessions. The system prevents overlapping availability entries and validates time slot formats. &
Pass \\ \midrule

F-11 & Scheduling -- Coordinator organizes session (UC-SCH-02) &
A coordinator opens the queue of sessions in state \textit{Proposed} or \textit{RescheduleRequested} via /scheduling/sessions endpoint, reviews the tutor's availability, assigns students, selects an appropriate time slot, chooses a room or online link, and finalizes the schedule. &
The session is updated with finalized time, location or online link, and a participant list containing one tutor and multiple students. The session state changes to \textit{Scheduled}. Notifications are sent to the tutor and all assigned students. &
Pass \\ \midrule

F-12 & Scheduling -- Request reschedule &
A student or tutor assigned to a \textit{Scheduled} session determines that the scheduled time conflicts with their personal schedule and submits a reschedule request via /scheduling/sessions/\{id\}/reschedule with a reason. &
The session time is updated or flagged for coordinator review. The session state may change to \textit{RescheduleRequested}. The coordinator is notified and the session reappears in their scheduling queue. &
Pass \\ \midrule

F-13 & Scheduling -- Manage student session (UC-SCH-03) &
A student logs in, opens the list of sessions via /sessions/my-sessions, views the details of a \textit{Scheduled} session, and can join the session or request changes if allowed. &
The system displays all sessions the student is assigned to with accurate schedules and participation status. The student can view session details, materials, and attendance records. &
Pass \\ \midrule

F-14 & Scheduling -- Change session time &
\textit{This test case was excluded from implementation as the feature is not currently used in the system workflow. Session time changes are handled through the reschedule request process (F-12).} &
\textit{N/A} &
N/A \\ \midrule

F-15 & Scheduling -- Manage tutor session (UC-SCH-04) &
A tutor logs in, opens their session list via /sessions/my-sessions (including \textit{Scheduled}, \textit{Proposed}, and \textit{RescheduleRequested} sessions), selects a session, and either confirms the schedule or requests rescheduling. &
The system displays all sessions the tutor is assigned to. The tutor can view session details, manage materials, mark attendance, and track session progress. &
Pass \\ \midrule

F-16 & Session Management -- Conduct session (UC-SES-01) &
At session time, the tutor opens the session details via /sessions/\{id\}/complete endpoint, marks the session as completed, records notes about the session content and learning outcomes, and saves. &
The session state changes to \textit{Completed}. Session completion timestamp and notes are stored and become available for reporting and progress tracking. &
Pass \\ \midrule

F-17 & Session Management -- Check attendance &
During or after a session, the tutor accesses the attendance interface via /sessions/\{id\}/attendance and marks each student as Present, Late, or Absent, with optional notes. &
Attendance records are saved with timestamps. Students can view their own attendance history. The data feeds into reports for participation tracking. &
Pass \\ \midrule

F-18 & Session Management -- Upload materials &
The tutor uploads learning materials (PDFs, slides, documents) to a \textit{Scheduled} or \textit{Completed} session via /sessions/\{id\}/materials endpoint. &
Files are successfully uploaded to storage, and download links are generated. Students assigned to the session can view and download these materials via /sessions/\{id\}/materials. File size and type validations are enforced. &
Pass \\ \midrule

F-19 & Session Management -- View materials &
A student opens the session detail page via /sessions/\{id\}/materials and views the list of uploaded materials for that session. &
All materials uploaded by the tutor are displayed with file names, upload dates, and download links. Only students assigned to the session can access these materials. &
Pass \\ \midrule

F-20 & Session Management -- Download materials &
A student clicks the download link for a learning material file via /sessions/\{id\}/materials/\{material\_id\}/download endpoint. &
The file downloads successfully. The system logs the download event (optional). If the file is stored externally (e.g., HCMUT\_LIBRARY), the system provides a valid access link. &
Pass \\ \midrule

F-21 & Session Management -- View learning progress (UC-SES-02) &
A student logs in and opens the Learning Progress page via /progress/students/\{id\}/progress. The system displays the list of completed sessions, attendance, tutor notes, and aggregated indicators such as number of sessions and progress trend. &
The student can only see data for sessions in which they are a participant, and cannot access other students' data. All information matches the recorded session and attendance data. &
Pass \\ \midrule

F-22 & Session Management -- Give feedback &
After a session is \textit{Completed}, a student accesses the feedback form via /sessions/\{id\}/feedback, rates the session quality (1-5 stars), provides written comments, and submits. &
Feedback is recorded in the database linked to the session. The tutor can view aggregated feedback ratings and comments. Feedback data is used in tutor performance reports. &
Pass \\ \midrule

F-23 & Report Module -- View Courses Report &
A coordinator or department chair opens the Report module via /reports/courses endpoint, filtering by department and semester. &
The system displays statistics including total sessions, participating students and tutors, average ratings, and completion rates for each course. Data is consistent with the underlying database. &
Pass \\ \midrule

F-24 & Admin Configuration -- Manage User Accounts &
An administrator accesses the User Management interface via /admin/users, views the list of all users, searches for a specific user by email or name, and edits user information or account status (active/inactive). &
User records are updated successfully. Changes are reflected immediately in the system. User authentication and authorization adapt to the new account status. &
Pass \\ \midrule

F-25 & Learning Forum -- Create Study Group &
A student or tutor creates a new study group for a course via /study-groups/ endpoint, specifies the name, description, subject, and initial settings. &
The study group is created and visible to eligible users. Members can discuss topics, share materials, and organize collaborative learning activities. Group data is persisted correctly. &
Pass \\
\end{longtable}
\end{center}

% --------------------------------------------------

% --------------------------------------------------
\subsection{Non-functional Test Cases}

% ---------- Non-functional longtable ----------
\refstepcounter{table}\label{tab:nonfunctional-testcases}
\begin{center}

\begin{longtable}
{p{0.06\textwidth}L{0.20\textwidth}L{0.30\textwidth}L{0.30\textwidth}p{0.05\textwidth}}
\toprule
\textbf{TC ID} & \textbf{Category} & \textbf{Scenario} & \textbf{Expected result} & \textbf{Result} \\
\midrule
\endfirsthead

\toprule
\textbf{TC ID} & \textbf{Category} & \textbf{Scenario} & \textbf{Expected result} & \textbf{Result} \\
\midrule
\endhead

\midrule
\multicolumn{5}{r}{\textit{Continued on next page}}\\
\bottomrule
\endfoot

\bottomrule
\caption*{\textbf{Table \thetable}. Non-functional test cases for performance, security, scalability, load, and database integrity.}
\endlastfoot

NF-01 & Performance &
Around 30 concurrent users log in using HCMUT\_SSO and authenticate through /auth/login endpoint simultaneously (including students, tutors, and coordinators). Load test measures P95 response time. &
At least 95\% of login requests complete within 15 seconds under concurrent load. No 5xx errors or timeouts occur. System maintains stability under authentication burst. &
Pass \\ \midrule

NF-02 & Performance &
A coordinator opens the dashboard via /users/stats/coordinator endpoint which aggregates statistics from multiple sessions, registrations, and user data. &
The dashboard statistics load within 5 seconds. The UI remains responsive during data aggregation. API returns comprehensive coordinator statistics without timeout. &
Pass \\ \midrule

NF-03 & Security &
A student attempts to access coordinator-only endpoints such as /coordinator/tutor-registrations via direct URL or API request. &
The system responds with 403 Forbidden or 404 Not Found. No sensitive coordinator data is leaked. Role-based access control is properly enforced. &
Pass \\ \midrule

NF-04 & Security &
A non-participant student attempts to access session materials via /sessions/\{id\}/materials for a session they are not assigned to. &
The system allows access but returns empty materials list or denies access appropriately. Authorization checks prevent unauthorized material access. &
Pass \\ \midrule

NF-05 & Scalability &
A coordinator queries a large dataset via /sessions/ endpoint with limit parameter set to 100, simulating a system with hundreds of sessions. &
List and filter operations complete within 4 seconds without query timeouts. Pagination works correctly. The system handles large result sets efficiently. &
Pass \\ \midrule

NF-06 & Scalability &
Multiple students concurrently attempt to join the same session via /sessions/\{id\}/join endpoint, simulating 10 simultaneous join requests. &
No data corruption occurs. All concurrent operations are processed correctly. Final participant list is consistent. Database transactions maintain data integrity. &
Pass \\ \midrule

NF-07 & Load \& Stability &
A load test simulates 50 concurrent users performing session list queries via /sessions/ for sustained period, maintaining active connections. &
The system remains stable, response times stay acceptable, and no crashes or memory leaks occur. Server handles sustained concurrent load without degradation. &
Pass \\ \midrule

NF-08 & Load \& Stability &
During a peak period, 70 students submit registrations concurrently via /students/register endpoint within a short time window. &
All valid registrations are processed correctly. Error rate remains low. The system handles registration burst without service degradation. Queue processing works as expected. &
Pass \\ \midrule

NF-09 & Database integrity &
A coordinator attempts to create sessions with overlapping times for the same tutor via /tutors/check-schedule-conflicts endpoint. &
The system detects schedule conflicts and prevents double-booking. Validation returns meaningful error messages. Schedule data remains consistent across all views. &
Pass \\ \midrule

NF-10 & Database integrity &
An admin attempts to delete a user account via /admin/users/\{id\} endpoint that has associated session data and historical records. &
System blocks hard deletion or applies soft-delete mechanism while preserving historical data. No orphan records appear. Referential integrity is maintained. &
Pass \\
\end{longtable}

\end{center}


\paragraph*{Testing Summary}


The test cases presented above represent the core validation scenarios executed during system development and validated against the production API endpoint (https://tutor-suporting-system-production.up.railway.app/api/v1). These functional and non-functional tests cover the critical paths through all major modules: Account Management, Registration, Scheduling, Session Management, Reporting, Administration, and Learning Forum.

A total of 34 test cases were executed: 24 functional test cases (F-01 to F-25, excluding F-14 which was deemed unnecessary for the current system workflow) and 10 non-functional test cases (NF-01 to NF-10). All 34 test cases achieved \textit{Pass} status, demonstrating that the system meets its specified requirements under both normal and stress conditions.

The test suite was implemented using pytest with async HTTP client (httpx.AsyncClient) to validate API endpoints, authentication flows, authorization controls, data persistence, and system behavior under concurrent load. Performance tests measured P95 response times, while security tests verified role-based access control and data protection mechanisms.

Beyond these representative cases, additional detailed testing was performed to validate edge cases, error handling, input validation, API integration, and UI/UX behaviors. The comprehensive test suite, maintained in the development repository at backend/testcases/test\_all\_35\_cases.py, provides confidence in the system's robustness and readiness for deployment.

\paragraph*{Test Exclusions and Rationale}

Test case F-14 (Change session time) was excluded from the final test suite as this functionality is adequately covered by the reschedule request workflow (F-12). In the current system architecture, schedule modifications are handled through the coordinator's reschedule review process rather than direct time changes, which provides better audit trails and notification workflows. This design decision was validated through user acceptance testing with the coordinator role.
